\section{Introducción}

El trabajo práctico final integra tres herramientas utilizadas durante la materia: Python para procesar datos experimentales, LTspice para simular circuitos y Altium Designer para implementar una placa. La consigna exige que los gráficos de respuesta en frecuencia se presenten en escala semilogarítmica, con frecuencia en hertz y magnitud en decibeles, y que las diferencias entre teoría, simulación y medición sean analizadas.

El informe se organiza en cuatro bloques principales. Primero se describe la interfaz gráfica desarrollada. Luego se estudia el circuito transitorio RLC de la consigna. A continuación se analizan los filtros de segundo orden propuestos. Finalmente se presenta el diseño de un filtro pasabanda de dos etapas para voz humana y su implementación en PCB.

\subsection{Datos de trabajo}

Se utilizó el valor de capacitor correspondiente al grupo 8 en las tablas de la consigna:

\begin{equation}
    C_{\mathrm{transitorio}}=\SI{47}{nF},
    \qquad
    C_{\mathrm{filtros}}=\SI{82}{nF}.
\end{equation}

Para los filtros de segundo orden se tomó:

\begin{equation}
    R=\SI{50}{\ohm}, \qquad L=\SI{1}{mH}, \qquad C=\SI{82}{nF}.
\end{equation}

A partir de estos valores se obtiene la frecuencia natural común de los circuitos RLC:

\begin{equation}
    \omega_0=\frac{1}{\sqrt{LC}},
    \qquad
    f_0=\frac{1}{2\pi\sqrt{LC}}.
\end{equation}

Reemplazando:

\begin{equation}
    f_0=\frac{1}{2\pi\sqrt{\SI{1}{mH}\cdot \SI{82}{nF}}}
    \approx \SI{17.6}{kHz}.
\end{equation}

El factor de calidad de las etapas serie es:

\begin{equation}
    Q=\frac{\omega_0 L}{R}=\frac{1}{R}\sqrt{\frac{L}{C}}
    \approx 2.21.
\end{equation}
